\documentclass[11pt,a4paper]{article}
\usepackage[T1]{fontenc}
\usepackage{isabelle,isabellesym}

\usepackage{amsmath}
\usepackage{amssymb}
\usepackage{stmaryrd}

% this should be the last package used
\usepackage{pdfsetup}

\urlstyle{rm}
\isabellestyle{it}

\begin{document}

\title{Semicontinuity: Envelopes, Berge's Theorem and Measurable Selection}
\author{Dmitriy Traytel}
\maketitle

\begin{abstract}
  Upper and lower semicontinuity as optimal control and viscosity-solution
  arguments use them, for real-valued functions on a metric space.

  The $\varepsilon$-$\delta$ predicates come with their elementary calculus
  --- closure under addition, positive scaling and subtracting a continuous
  function --- and with the two facts such arguments actually consume:
  attainment of the supremum or infimum on a nonempty compact set, and
  bounded extension of an upper semicontinuous function off a closed set.
  The semicontinuous envelopes are taken both over the whole space and
  \emph{relative to a closed set}~$K$, the latter with a projection-based
  extension to the ambient space; the relative form is what a boundary
  condition stated in the viscosity sense actually reads.  Berge's maximum
  theorem is proved in the form that needs no measurable selection: a
  supremum of a jointly upper semicontinuous function over a compact index
  set is upper semicontinuous.  Finally, a measurable selection theorem for
  upper semicontinuous payoffs on compact sets, Bertsekas--Shreve
  Proposition~7.33~\cite{BertsekasShreve}, proved by a greedy nested
  bisection along a countable dense sequence; a countably valued
  $\varepsilon$-selector cannot exist, so some such construction is forced.

  The envelope defined here is the lower semicontinuous hull, which the
  Archive entry \texttt{Lower\_Semicontinuous} already provides over
  $\overline{\mathbb{R}}$ as \texttt{lsc\_hull}.  Rather than repeat it,
  this entry proves the two agree --- the hull's \emph{liminf} form is
  exactly the supremum of infima over balls used here --- and inherits the
  four properties that follow: the envelope lies below the function, above
  any lower bound, is monotone, and is itself lower semicontinuous.  The
  boundedness hypotheses carried throughout are what keep the values finite,
  which is the price of working in $\mathbb{R}$ rather than
  $\overline{\mathbb{R}}$; that entry's own subject, convex functions and
  their epigraphs, is disjoint from this one's.

  Extracted from a formalisation of the uniqueness half of a
  relative-arbitrage problem~\cite{LaiShkolnikovSoner}, whose value function
  is characterised as the unique bounded upper semicontinuous viscosity
  solution of a degenerate elliptic equation.
\end{abstract}

\tableofcontents

\parindent 0pt\parskip 0.5ex

\input{session}

\bibliographystyle{abbrv}
\bibliography{root}

\end{document}
